
%%%%%%%%%%%%%%%%%%%%%%%%%%%%%%%%%%%%%%%%%%%%%%%%%%%%%%%%%%%%%
\section{模型的评价}

\subsection{模型的优点}
\begin{itemize}[itemindent=2em]
\item \textbf{机理为主、外推可靠：}以Deutsch方程等物理机理模型为主框架，从历史数据只拟合待定系数，外推时靠物理单调性约束保证趋势合理，避免了纯数据驱动在$C_{out}$方差极小时外推不可靠的问题。
\item \textbf{工况自适应寻优：}K-Means工况划分使不同入口条件下的最优参数独立求解，避免了"一刀切"统一参数的弊端，高浓度与低浓度工况策略差异显著，符合工程实际。
\item \textbf{排放阈值参数化设计：}问题四复用问题二的优化框架，仅需传入不同$C_{limit}$即可求解，代码与模型结构一致性好，便于扩展到其他排放标准。
\item \textbf{多起点+全局备选保证收敛：}SLSQP多起点启动+差分进化备选，有效应对非凸可行域，6个工况均成功收敛且约束满足。
\end{itemize}

\subsection{模型的缺点}
\begin{itemize}[itemindent=2em]
\item \textbf{振打模型单向局限：}当前仅建模"周期过长效率下降"，未建模"周期过短清灰更彻底效率提升"的双向效应，导致振打弹性性价比为零，优化器总将振打压到下界。
\item \textbf{振打峰值估算粗糙：}半经验公式中系数0.1为经验值，缺乏标定数据；分钟级数据无法直接观测秒级峰值，估算不确定性较大。
\item \textbf{电场间耦合未建模：}假设$H_1$认为各级独立，但前电场振打释放的粉尘会重返气流进入后电场，振打瞬间级间独立性不成立。
\item \textbf{$C_{out}$外推不确定性：}从历史$C_{out}\approx 50$外推到$10$甚至$5\;\text{mg/Nm}^3$跨两个数量级，Deutsch参数拟合虽有物理约束，但数据欠定下参数不确定性仍较大。
\end{itemize}
